\paragrafos{Απόλυτη τιμή και διάταξη: ποιος είναι δεξιότερα;}{Α.3.2, «Απόλυτη τιμή ρητού, σύγκριση και διάταξη ρητών, αντίθετοι ρητοί»}

\begin{anagke}
Στην ώρα των μαθηματικών η κυρία Αλεξάνδρα ρωτάει: «Ποιος είναι ο επόμενος ρητός μετά το $\frac{2}{5}$;». Ο Βασίλης απαντάει με σιγουριά: «Το $\frac{3}{5}$». Η Ρήνη σηκώνει το χέρι: «Στον χάρακα το $\frac{2}{5}$ είναι στο $0{,}4$ και το $\frac{3}{5}$ στο $0{,}6$. Ανάμεσα υπάρχει το $0{,}5$. Και το $0{,}45$. Και το $0{,}41$!». Ο Νίκος ρωτάει: «Αν δεν υπάρχει επόμενος, πώς τους βάζουμε σε σειρά;».

Λίγο αργότερα ο Ηλίας διαβάζει την πρόγνωση: Ιωάννινα $-3^\circ$, Φλώρινα $-10^\circ$. «Το δέκα είναι μεγαλύτερο από το τρία, άρα η Φλώρινα έχει μεγαλύτερη θερμοκρασία», λέει ο Βασίλης. Η Ζωή διαφωνεί: «Δεν μπορεί να είναι πιο ζεστό το δέκα κάτω από το μηδέν!». Η Αγγελική ζωγραφίζει τον άξονα: «Όσο πιο αριστερά, τόσο πιο κρύο. Το δέκα μετράει πόσο μακριά είμαστε από το μηδέν, και το πρόσημο λέει από ποια μεριά».

Η Πολυξένη το διατυπώνει με δύο ερωτήματα: «Πόσο μακριά είναι κάθε αριθμός από το μηδέν; Και ποιος είναι πιο δεξιά;» Η πρώτη ιδέα θα λέγεται απόλυτη τιμή και η δεύτερη διάταξη. Σε αυτή την παράγραφο θα μάθουμε να συγκρίνουμε οποιουσδήποτε δύο ρητούς: κλάσματα, δεκαδικούς, θετικούς και αρνητικούς.
\end{anagke}

\begin{lexi}{απόλυτη τιμή, αντίθετοι, διάταξη}{απόλυτη τιμή, αντίθετοι, διάταξη}
\lexiword{Απόλυτη} σημαίνει «λυμένη από δεσμούς» (από το «απολύω»): η απόλυτη τιμή είναι ο αριθμός ελεύθερος από το πρόσημό του, δηλαδή μόνο η απόσταση από το $0$. Οι \lexiword{αντίθετοι} βρίσκονται σε «αντίθετη» κατεύθυνση: ίδια απόσταση από το $0$, ο ένας δεξιά και ο άλλος αριστερά. Η \lexiword{διάταξη} είναι η «διά-ταξη», το να βάζεις τους αριθμούς «σε τάξη», δηλαδή σε σειρά.
\end{lexi}

\begin{metafrasi}
Η απόσταση του $-2$ από το $0$ είναι $2$. & $|-2|=2$\\
Η απόσταση του $+\frac{1}{2}$ από το $0$ είναι $\frac{1}{2}$. & $\left|+\frac{1}{2}\right|=\frac{1}{2}$\\
Ο αντίθετος του $+2{,}8$ είναι ο $-2{,}8$. & $-(+2{,}8)=-2{,}8$\\
Το $-3{,}7$ είναι αριστερά του $-2{,}5$. & $-3{,}7<-2{,}5$\\
Το $1{,}6$ είναι δεξιά του $0{,}5$. & $1{,}6>0{,}5$\\
\end{metafrasi}

\tmima{Η θεωρία}

\tmimasmall{Απόλυτη τιμή και αντίθετοι}

Ο ορισμός της απόλυτης τιμής που ξέρουμε από τους ακεραίους ισχύει και για τους ρητούς.

\begin{orismos}{απόλυτη τιμή ρητού}
\textbf{Απόλυτη τιμή} ενός ρητού $\alpha$ είναι η \textbf{απόσταση} του σημείου με τετμημένη $\alpha$ από την αρχή $\mathrm{Ο}$ του άξονα. Συμβολίζεται $|\alpha|$.
\end{orismos}

Για παράδειγμα $|-2|=2$ και $\left|+\frac{1}{2}\right|=\frac{1}{2}$. Η απόλυτη τιμή ενός θετικού είναι ο ίδιος ο αριθμός, ενός αρνητικού ο αντίθετός του και $|0|=0$. Δεν είναι ποτέ αρνητική.

\begin{center}
\akgr[1.8]{-2}{2}{-2/-2,0.5/\frac{1}{2}}
\end{center}

\begin{orismos}{αντίθετοι ρητοί}
\textbf{Αντίθετοι} λέγονται δύο ρητοί που είναι \textbf{ετερόσημοι} και έχουν την \textbf{ίδια απόλυτη τιμή}. Ο αντίθετος του $\alpha$ είναι ο $-\alpha$. Ο αντίθετος του $0$ είναι το $0$ και ισχύει $-(-\alpha)=\alpha$.
\end{orismos}

Οι $+2{,}8$ και $-2{,}8$ είναι αντίθετοι. Το ίδιο και οι $-\frac{1}{7}$ και $+\frac{1}{7}$. Στον άξονα, οι αντίθετοι είναι σημεία συμμετρικά ως προς το $0$.

\tmimasmall{Σύγκριση και διάταξη ρητών}

Από δύο ρητούς \textbf{μεγαλύτερος} είναι αυτός που βρίσκεται \textbf{δεξιότερα} στον άξονα.

\begin{center}
\akgr[1.8]{-4}{2}{-3.7/-3{,}7,-2.5/-2{,}5,0.5/0{,}5,1.6/1{,}6}
\end{center}

Στο σχήμα $1{,}6>0{,}5$ και $-2{,}5>-3{,}7$.

\begin{kanonas}{σύγκριση ρητών με απόλυτες τιμές}
\begin{itemize}[leftmargin=1.2em,itemsep=0.3ex,topsep=2pt]
\item Κάθε θετικός είναι μεγαλύτερος από κάθε αρνητικό. Το $0$ είναι μεγαλύτερο από κάθε αρνητικό και μικρότερο από κάθε θετικό.
\item Από δύο \textbf{θετικούς} μεγαλύτερος είναι αυτός που έχει τη \textbf{μεγαλύτερη} απόλυτη τιμή: $+4{,}3>+4{,}2$, γιατί $|{+4{,}3}|>|{+4{,}2}|$.
\item Από δύο \textbf{αρνητικούς} μεγαλύτερος είναι αυτός που έχει τη \textbf{μικρότερη} απόλυτη τιμή: $-0{,}1>-0{,}2$, γιατί $|{-0{,}1}|=0{,}1<0{,}2=|{-0{,}2}|$.
\end{itemize}
\end{kanonas}

\tmimasmall{Σύγκριση κλασμάτων}

Για κλάσματα με \textbf{θετικούς} όρους ισχύουν τα εξής.

\begin{kanonas}{σύγκριση κλασμάτων}
\begin{itemize}[leftmargin=1.2em,itemsep=0.3ex,topsep=2pt]
\item Από δύο \textbf{ομώνυμα} κλάσματα μεγαλύτερο είναι αυτό με τον \textbf{μεγαλύτερο αριθμητή}: $\frac{2}{5}<\frac{3}{5}$.
\item Από δύο κλάσματα με τον \textbf{ίδιο αριθμητή} μεγαλύτερο είναι αυτό με τον \textbf{μικρότερο παρονομαστή}: $\frac{1}{4}<\frac{1}{2}$.
\item Για να συγκρίνουμε \textbf{ετερώνυμα} κλάσματα, τα κάνουμε \textbf{ομώνυμα} και συγκρίνουμε τους αριθμητές τους. Διαφορετικά, εκτελούμε τις διαιρέσεις και συγκρίνουμε τους δεκαδικούς.
\end{itemize}
\end{kanonas}

\begin{center}
\(\dfrac{2}{5}\)\ \binpie{5}{2}\ \ $<$\ \ \(\dfrac{3}{5}\)\ \binpie{5}{3}\qquad\qquad\(\dfrac{1}{4}\)\ \binpie{4}{1}\ \ $<$\ \ \(\dfrac{1}{2}\)\ \binpie{2}{1}
\end{center}

\tmimasmall{Σύγκριση δεκαδικών}

\begin{kanonas}{σύγκριση δεκαδικών}
\begin{itemize}[leftmargin=1.2em,itemsep=0.3ex,topsep=2pt]
\item Αν έχουν διαφορετικό ακέραιο μέρος, μεγαλύτερος είναι αυτός με το μεγαλύτερο ακέραιο μέρος: $8{,}43>5{,}96$.
\item Αν έχουν το ίδιο ακέραιο μέρος, συγκρίνουμε τα δεκαδικά τους ψηφία \textbf{ένα προς ένα από αριστερά προς τα δεξιά} και μεγαλύτερος είναι αυτός που έχει μεγαλύτερο το πρώτο ψηφίο στο οποίο διαφέρουν: $3{,}132>3{,}1294$, γιατί στα χιλιοστά $3>2$.
\end{itemize}
\end{kanonas}

\tmimasmall{Άπειροι ρητοί ανάμεσα σε δύο ρητούς}

Ανάμεσα σε δύο οποιουσδήποτε διαφορετικούς ρητούς υπάρχει ένα \textbf{άπειρο πλήθος} ρητών. Λέμε ότι οι ρητοί είναι «\textbf{πυκνοί}» στην ευθεία: δεν υπάρχει «επόμενος» ενός ρητού, σε αντίθεση με τους ακεραίους, όπου κάθε ακέραιος έχει επόμενο. Για να βρούμε έναν ρητό ανάμεσα σε δύο άλλους, τους γράφουμε ως ομώνυμα (και αν δεν υπάρχει ακέραιος ανάμεσα στους αριθμητές, πολλαπλασιάζουμε τους όρους με έναν αρκετά μεγάλο αριθμό, π.χ. $10$) ή ως δεκαδικούς.

\begin{prosoxi}
\begin{itemize}[leftmargin=1.2em,itemsep=0.3ex,topsep=2pt]
\item Στους αρνητικούς το «μεγάλο νούμερο» είναι μικρός αριθμός: $-10<-3$ παρά το $10>3$.
\item Ο αντίθετος του $-5$ είναι το $+5$ και όχι το $-5$. Και $|-5|=5$, όχι $-5$.
\item Στα δεκαδικά, ο αριθμός με τα περισσότερα ψηφία δεν είναι μεγαλύτερος: $0{,}3>0{,}2999$. Και ο $-1{,}36$ είναι \emph{μικρότερος} από τον $-1{,}3$.
\item Στα κλάσματα, ο μεγαλύτερος παρονομαστής δεν δίνει μεγαλύτερο κλάσμα: $\frac{1}{8}<\frac{1}{3}$. Τα κομμάτια είναι πιο μικρά.
\item Δεν συγκρίνουμε ετερώνυμα κλάσματα κοιτάζοντας μόνο τους αριθμητές ή μόνο τους παρονομαστές.
\end{itemize}
\end{prosoxi}

\tmima{Λυμένα παραδείγματα}

\paradeigma{Ποια πόλη είναι πιο κρύα}
\tixero Ιωάννινα $-3^\circ$, Φλώρινα $-10^\circ$.\par
\tithelo Ποια θερμοκρασία είναι μικρότερη;\par
\possindeo Και οι δύο είναι αρνητικοί, και μεγαλύτερος είναι αυτός με τη μικρότερη απόλυτη τιμή. Ισχύει $|-10|=10>3=|-3|$.\par
\apantisi $-10<-3$: η Φλώρινα είναι πιο κρύα.\par
\elegxos Στον άξονα το $-10$ είναι αριστερότερα από το $-3$.

\paradeigma{Αντίθετοι και απόλυτες τιμές}
\tixero Οι ρητοί $+4{,}5$, $-\frac{2}{3}$ και $-7$.\par
\tithelo Να βρω τις απόλυτες τιμές και τους αντίθετούς τους.\par
\possindeo Η απόλυτη τιμή είναι η απόσταση από το $0$: $4{,}5$, $\frac{2}{3}$, $7$. Ο αντίθετος του $x$ είναι ο $-x$: $-4{,}5$, $+\frac{2}{3}$, $+7$.\par
\apantisi $|{+4{,}5}|=4{,}5$, $\left|-\frac{2}{3}\right|=\frac{2}{3}$, $|-7|=7$ και αντίθετοι $-4{,}5$, $+\frac{2}{3}$, $+7$.\par
\elegxos Κάθε αριθμός και ο αντίθετός του έχουν την ίδια απόλυτη τιμή.

\paradeigma{Ετερώνυμα κλάσματα}
\tixero Τα κλάσματα $\frac{3}{5}$, $\frac{5}{8}$ και $\frac{7}{12}$.\par
\tithelo Να τα βάλω σε αύξουσα σειρά.\par
\possindeo Τα κάνω ομώνυμα. $\EKP(5,8,12)=120$. $\frac{3}{5}=\frac{72}{120}$, $\frac{5}{8}=\frac{75}{120}$, $\frac{7}{12}=\frac{70}{120}$. Συγκρίνω τους αριθμητές: $70<72<75$.\par
\apantisi $\dfrac{7}{12}<\dfrac{3}{5}<\dfrac{5}{8}$.\par
\elegxos Με δεκαδικούς: $0{,}583\ldots<0{,}6<0{,}625$.

\paradeigma{Αρνητικά κλάσματα και δεκαδικοί}
\tixero Οι αριθμοί $-\frac{5}{2}$, $-0{,}9$, $-4{,}5$, $0{,}1$, $\frac{1}{2}$, $3$.\par
\tithelo Να τους τοποθετήσω στην αριθμογραμμή και να τους συγκρίνω.\par
\possindeo Γράφω τα κλάσματα ως δεκαδικούς: $-2{,}5$ και $0{,}5$. Οι αρνητικοί είναι αριστερά του $0$ και ο μικρότερος είναι αυτός με τη μεγαλύτερη απόλυτη τιμή: $-4{,}5<-2{,}5<-0{,}9$. Μετά οι θετικοί: $0{,}1<0{,}5<3$.\par
\apantisi $-4{,}5<-\dfrac{5}{2}<-0{,}9<0{,}1<\dfrac{1}{2}<3$.\par
\elegxos Η σειρά συμφωνεί με τη σειρά των σημείων στον άξονα, από αριστερά προς τα δεξιά.

\begin{center}
\akgr[1.1]{-5}{4}{-4.5/-4{,}5,-2.5/-\frac{5}{2},-0.9/-0{,}9,0.1/0{,}1,0.5/\frac{1}{2},3/3}
\end{center}

\paradeigma{Ένας ρητός ανάμεσα}
\tixero Οι αριθμοί $\frac{1}{3}$ και $\frac{1}{2}$.\par
\tithelo Να βρω έναν ρητό ανάμεσά τους.\par
\possindeo \emph{1ος τρόπος.} Κάνω τα κλάσματα ομώνυμα: $\frac{1}{3}=\frac{2}{6}$ και $\frac{1}{2}=\frac{3}{6}$. Δεν υπάρχει ακέραιος ανάμεσα στο $2$ και στο $3$, γι' αυτό πολλαπλασιάζω τους όρους επί $10$: $\frac{20}{60}$ και $\frac{30}{60}$. Ανάμεσα βρίσκεται το $\frac{25}{60}$. \emph{2ος τρόπος.} Με δεκαδικούς: $\frac{1}{3}=0{,}333\ldots$ και $\frac{1}{2}=0{,}5$, άρα ανάμεσα βρίσκεται ο $0{,}4$.\par
\apantisi $\dfrac{1}{3}<\dfrac{25}{60}<\dfrac{1}{2}$ και $\dfrac{1}{3}<0{,}4<\dfrac{1}{2}$.\par
\elegxos $\frac{25}{60}=\frac{5}{12}=0{,}41\overline{6}$. Βρίσκεται ανάμεσα στο $0{,}333\ldots$ και στο $0{,}5$.

\paradeigma{Σύγκριση δεκαδικών}
\tixero Οι αριθμοί $9{,}8$ και $9{,}75$, και οι $-1{,}36$ και $-1{,}3$.\par
\tithelo Ποιος είναι μεγαλύτερος σε κάθε ζεύγος;\par
\possindeo Ίδιο ακέραιο μέρος στο πρώτο ζεύγος. Γράφω $9{,}8=9{,}80$ και συγκρίνω τα δέκατα: $8>7$, άρα $9{,}80>9{,}75$. Στο δεύτερο οι αριθμοί είναι αρνητικοί, και $|-1{,}36|=1{,}36>1{,}30=|-1{,}3|$, άρα $-1{,}36<-1{,}3$.\par
\apantisi $9{,}8>9{,}75$ και $-1{,}3>-1{,}36$.\par
\elegxos Αν $9{,}8$ και $9{,}75$ είναι χρόνοι τρεξίματος, ο $9{,}75$ είναι ο καλύτερος.

\tmima{Ασκήσεις}

\zoni{Το κατάλαβα}{απόλυτη τιμή, αντίθετοι, σύγκριση}

\ex Συμπλήρωσε τα κενά. α) Ο αντίθετος του $\alpha$ είναι ο \ldots \quad β) Από δύο θετικούς ρητούς μεγαλύτερος είναι αυτός με τη \ldots απόλυτη τιμή. \quad γ) Από δύο αρνητικούς ρητούς μεγαλύτερος είναι αυτός με τη \ldots απόλυτη τιμή. \quad δ) Από δύο ομώνυμα κλάσματα μεγαλύτερο είναι αυτό με τον μεγαλύτερο \ldots \quad ε) Από δύο κλάσματα με ίδιο αριθμητή μεγαλύτερο είναι αυτό με τον \ldots παρονομαστή.

\ex Συμπλήρωσε τον πίνακα.

\begin{center}
\renewcommand{\arraystretch}{1.6}
\begin{tabular}{|l|c|c|c|c|c|c|}
\hline
Αριθμός & $-3$ & $0{,}35$ & $\frac{1}{9}$ & & & $-\frac{2}{5}$\\
\hline
Αντίθετος & & & & $-0{,}1$ & $\frac{7}{4}$ & \\
\hline
Απόλυτη τιμή & & & & & & \\
\hline
\end{tabular}
\end{center}

\ex Σύγκρινε με ένα από τα $<,\ >,\ =$: α) $\frac{21}{6}\ \ldots\ \frac{12}{6}$ \quad β) $\frac{4}{9}\ \ldots\ \frac{4}{6}$ \quad γ) $\frac{2}{3}\ \ldots\ \frac{7}{8}$ \quad δ) $\frac{1}{100}\ \ldots\ -\frac{2}{5}$ \quad ε) $-\frac{8}{10}\ \ldots\ -1$ \quad στ) $-\frac{8}{3}\ \ldots\ -3$ \quad ζ) $\frac{9}{8}\ \ldots\ 1{,}2$.

\ex Σύγκρινε τους δεκαδικούς: α) $3{,}16\ \ldots\ 3{,}2$ \quad β) $0{,}092\ \ldots\ 0{,}087$ \quad γ) $-1{,}36\ \ldots\ -0{,}01$ \quad δ) $-2{,}69\ \ldots\ 1{,}4$ \quad ε) $-1{,}3\ \ldots\ -1{,}33$ \quad στ) $-1{,}10\ \ldots\ -1{,}1$.

\ex Βάλε σε αύξουσα σειρά τους δεκαδικούς: $1,\ -3,\ 0{,}2,\ 1{,}37,\ -0{,}03,\ 0$.

\ex Βάλε σε αύξουσα σειρά τα κλάσματα: α) $\frac{5}{4},\ \frac{1}{4},\ -\frac{3}{4},\ \frac{3}{4},\ -\frac{1}{4},\ -\frac{5}{4}$ \quad β) $\frac{5}{3},\ \frac{5}{4},\ -\frac{5}{2},\ \frac{5}{8},\ -\frac{5}{7},\ -\frac{5}{3}$.

\ex Δίνονται τα κλάσματα $\frac{9}{4},\ \frac{3}{16},\ \frac{9}{8},\ \frac{15}{16},\ \frac{5}{4},\ \frac{7}{16}$. α) Βρες ποια είναι μικρότερα του $1$ και γράψε τα σε αύξουσα σειρά. β) Βρες ποια είναι μεγαλύτερα του $1$ και γράψε τα σε αύξουσα σειρά. γ) Γράψε και τα έξι σε αύξουσα σειρά.

\ex Σωστό ή λάθος; α) $|\alpha|\ge0$ για κάθε ρητό $\alpha$. \quad β) $|-5|=-5$. \quad γ) Ο αντίθετος του $-4$ είναι ο $4$. \quad δ) $-9>-2$. \quad ε) $\frac{3}{7}<\frac{3}{5}$. \quad στ) Ο αντίθετος του $0$ είναι το $0$.

\zoni{Το χρησιμοποιώ}{θερμόμετρα, ώρες, αθλήματα}

\ex \textbf{Οι θερμοκρασίες.} Στον πίνακα του καιρού: Μέτσοβο $-5{,}5^\circ$, Ιωάννινα $-3^\circ$, Φλώρινα $-10^\circ$, Αθήνα $7^\circ$, Θεσσαλονίκη $-1{,}5^\circ$. Γράψε τις πόλεις από την πιο κρύα στην πιο ζεστή.

\ex \textbf{Οι δραστηριότητες.} Στο κυκλικό διάγραμμα της Αγγελικής ο Ύπνος παίρνει το $\frac{1}{3}$ του εικοσιτετραώρου, το Σχολείο το $\frac{1}{4}$, το Διάβασμα το $\frac{1}{8}$ και η Διασκέδαση το $\frac{1}{6}$. α) Πόσες ώρες διαρκεί η κάθε δραστηριότητα; β) Γράψε τις δραστηριότητες από αυτή που παίρνει τον περισσότερο χρόνο ως αυτή που παίρνει τον λιγότερο. γ) Τι μέρος του εικοσιτετραώρου μένει για όλα τα άλλα;

\ex \textbf{Τα εκατοστά του δευτερολέπτου.} Ο Νίκος τρέχει τα $60$ μέτρα σε $9{,}8$ δευτερόλεπτα και ο Βασίλης σε $9{,}75$. Ο Βασίλης λέει «το $75$ είναι μεγαλύτερο από το $8$, άρα έκανα περισσότερο χρόνο». Ποιος κέρδισε και ποιο είναι το λάθος του Βασίλη;

\ex \textbf{Οι αποστάσεις.} Στον άξονα της οδού το σχολείο είναι στο $0$ (μονάδα $100$ μ.), ο φούρνος στο $+3$ και το βιβλιοπωλείο στο $-2$. Ποιο μαγαζί είναι πιο κοντά στο σχολείο; Πόσο απέχουν μεταξύ τους τα δύο μαγαζιά;

\ex \textbf{Τα υψόμετρα.} Η Νεκρά Θάλασσα είναι στα $-430$ μ., ο Όλυμπος στα $+2\,917$ μ., η επιφάνεια της θάλασσας στο $0$ και ο βυθός της τάφρου των Μαριανών στα $-11\,034$ μ. Γράψε τα σε αύξουσα σειρά υψομέτρου.

\zoni{Το σκέφτομαι}{γιατί δεν υπάρχει επόμενος;}

\ex Ο Βασίλης λέει ότι ο επόμενος ρητός του $\frac{2}{5}$ είναι ο $\frac{3}{5}$. Βρες τρεις ρητούς ανάμεσά τους.

\ex Βρες δύο ρητούς ανάμεσα στους: α) $-0{,}7$ και $-0{,}8$ \quad β) $\frac{3}{10}$ και $\frac{2}{5}$ \quad γ) $\frac{7}{9}$ και $\frac{5}{6}$.

\ex Για ποιους ρητούς ισχύει $|\alpha|=\alpha$ και για ποιους $|\alpha|=-\alpha$; Γιατί ισχύει πάντα $|\alpha|=|-\alpha|$;

\ex Βρες τους δύο διαδοχικούς ακεραίους ανάμεσα στους οποίους βρίσκεται κάθε αριθμός: α) $0{,}3$ \quad β) $-1{,}25$ \quad γ) $\frac{7}{2}$ \quad δ) $-\frac{11}{3}$ \quad ε) $\frac{2024}{2025}$.

\zoni{Για το σπίτι}{ψάξε και σύγκρινε}

\ex Βρες τις θερμοκρασίες πέντε πόλεων του κόσμου και βάλε τις σε σειρά από την πιο κρύα στην πιο ζεστή. Σχεδίασε και τον άξονα.

\ex Ζήτησε από τρία μέλη της οικογένειάς σου να σου πουν τον χρόνο που κοιμούνται (σε ώρες) και γράψε το κλάσμα του εικοσιτετραώρου. Ποιος κοιμάται περισσότερο;
